\documentclass{article}
\usepackage{amsmath,amssymb,amsthm,algorithm,algorithmic}
\usepackage{hyperref}
\usepackage{bm}
\usepackage{enumitem}
\newtheorem{theorem}{Theorem}
\newtheorem{lemma}{Lemma}
\newtheorem{assumption}{Assumption}
\newtheorem{corollary}{Corollary}
\newtheorem{definition}{Definition}
\begin{document}

%%%%%%%%%%%%%%%%%%%%%%%%%%%%%%%%%%%%%%%%%%%%%%%%%%%%%%%%%%%%
\section*{Algorithm Name}
\textbf{Stochastic Weight Averaging with Interpolated Snapshots (SWA‑I)}  

The method keeps exact copies of the iterate every $k$ steps (snapshots) and
forms a simple linear interpolation between two successive snapshots for
prediction.  The snapshots are later averaged to obtain a final
model.

%%%%%%%%%%%%%%%%%%%%%%%%%%%%%%%%%%%%%%%%%%%%%%%%%%%%%%%%%%%%
\section*{Mathematical Formulation}
Let $f:\mathbb{R}^{d}\rightarrow\mathbb{R}$ be the (possibly non‑convex) loss
function and let $\nabla f(w)$ denote its gradient.  
For a stepsize $\eta>0$ and a snapshot period $k\in\mathbb{N}$ we define
the iterates
\[
\boxed{%
w_{t+1}=w_{t}-\eta\,\nabla f(w_{t}),\qquad t=0,1,\dots,T-1.}
\tag{1}
\]

\noindent
\textbf{Snapshots.}
Define the set of snapshot indices
\[
\mathcal{S}:=\{\,i k\mid i=1,2,\dots,M\,\},\qquad M:=\Big\lfloor\frac{T}{k}\Big\rfloor .
\tag{2}
\]
The corresponding snapshot vectors are $w_{i k}$, $i=1,\dots,M$.

\noindent
\textbf{Linear interpolation.}
For any $t\in[(i-1)k,i k]$ we write
\[
\alpha_{t}:=\frac{t-(i-1)k}{k}\in[0,1],\qquad
\tilde w_{t}:=(1-\alpha_{t})\,w_{(i-1)k}+\alpha_{t}\,w_{i k}.
\tag{3}
\]
The interpolated point $\tilde w_{t}$ may be used for prediction or
evaluation during training.

\noindent
\textbf{Final averaged model.}
After $T$ steps the output of SWA‑I is the simple average of all snapshots
\[
\boxed{\;\bar w_{T}:=\frac{1}{M}\sum_{i=1}^{M} w_{i k}\; }.
\tag{4}
\]

%%%%%%%%%%%%%%%%%%%%%%%%%%%%%%%%%%%%%%%%%%%%%%%%%%%%%%%%%%%%
\section*{Pseudocode}
\begin{algorithm}[H]
\caption{SWA‑I}
\begin{algorithmic}[1]
\STATE \textbf{Input:} initial point $w_{0}\in\mathbb{R}^{d}$, stepsize $\eta$, snapshot period $k$, total iterations $T$.
\STATE Initialize $w\leftarrow w_{0}$, $M\leftarrow 0$, $\mathcal{W}\leftarrow\emptyset$.
\FOR{$t=0$ \TO $T-1$}
    \STATE Compute stochastic gradient $g_{t}=\nabla f(w)$   \hfill \textit{(or mini‑batch estimate)}
    \STATE $w\leftarrow w-\eta g_{t}$ \hfill \textit{(gradient step (1))}
    \IF{$(t+1)\bmod k =0$}
        \STATE $M\leftarrow M+1$;   $\mathcal{W}\leftarrow\mathcal{W}\cup\{w\}$   \hfill \textit{store snapshot}
    \ENDIF
\ENDFOR
\STATE $\displaystyle \bar w_{T}\;:=\;\frac{1}{M}\sum_{w\in\mathcal{W}} w$ \hfill \textit{average (4)}
\RETURN $\bar w_{T}$
\end{algorithmic}
\end{algorithm}

%%%%%%%%%%%%%%%%%%%%%%%%%%%%%%%%%%%%%%%%%%%%%%%%%%%%%%%%%%%%
\section*{Dry Run Example}
Consider the simple quadratic
\[
f(w)=(w-3)^{2},\qquad \nabla f(w)=2(w-3).
\]
Set $w_{0}=0$, stepsize $\eta=0.1$, snapshot period $k=2$, and run three
gradient steps ($T=3$).

\begin{center}
\begin{tabular}{c|c|c|c}
$t$ & $w_{t}$ & $\nabla f(w_{t})$ & $f(w_{t})$ \\ \hline
0 & $0$ & $2(0-3)=-6$ & $(0-3)^{2}=9$ \\ 
1 & $w_{1}=0-0.1(-6)=0.6$ & $2(0.6-3)=-4.8$ & $(0.6-3)^{2}=5.76$ \\ 
2 & $w_{2}=0.6-0.1(-4.8)=1.08$ & $2(1.08-3)=-3.84$ & $(1.08-3)^{2}=3.6864$ \\ 
3 & $w_{3}=1.08-0.1(-3.84)=1.464$ & $2(1.464-3)=-3.072$ & $(1.464-3)^{2}=2.3689$
\end{tabular}
\end{center}

Snapshots are taken at $t=2$ (since $k=2$), i.e. $w_{2}=1.08$.  
The final averaged model after $T=3$ steps would be $\bar w_{3}=w_{2}=1.08$
(because only one snapshot exists).  
Interpolation for $t=1$ (between $w_{0}$ and $w_{2}$) gives
\[
\alpha_{1}= \frac{1-0}{2}=0.5,\qquad
\tilde w_{1}=0.5\,w_{0}+0.5\,w_{2}=0.5\times0+0.5\times1.08=0.54,
\]
which is close to the actual $w_{1}=0.6$.

%%%%%%%%%%%%%%%%%%%%%%%%%%%%%%%%%%%%%%%%%%%%%%%%%%%%%%%%%%%%
\section*{Assumptions}
\begin{assumption}[Smoothness]\label{as:smooth}
$f$ is $L$‑smooth, i.e.
\[
\|\nabla f(x)-\nabla f(y)\|\le L\|x-y\|\qquad\forall x,y\in\mathbb{R}^{d}.
\tag{A1}
\]
Equivalently,
\[
f(y)\le f(x)+\langle\nabla f(x),y-x\rangle+\frac{L}{2}\|y-x\|^{2}.
\tag{A1'}
\]
\end{assumption}

\begin{assumption}[Bounded Below]\label{as:lower}
$f$ is bounded from below: $f^{\star}:=\inf_{w}f(w)>-\infty$.
\tag{A2}
\end{assumption}

\begin{assumption}[Stepsize]\label{as:step}
The constant stepsize satisfies
\[
0<\eta\le\frac{1}{L}.
\tag{A3}
\]
\end{assumption}

No convexity is required; the analysis holds for general non‑convex smooth
functions.

%%%%%%%%%%%%%%%%%%%%%%%%%%%%%%%%%%%%%%%%%%%%%%%%%%%%%%%%%%%%
\section*{Main Convergence Theorem}
\begin{theorem}[Convergence of SWA‑I]\label{thm:main}
Under Assumptions \ref{as:smooth}–\ref{as:step}, let $\{w_{t}\}_{t=0}^{T}$
be generated by \eqref{1} and let $\{w_{i k}\}_{i=1}^{M}$ be the
snapshots defined in \eqref{2}.  Then the average squared gradient norm
over the snapshot iterates satisfies
\[
\boxed{\;
\frac{1}{M}\sum_{i=1}^{M}\bigl\|\nabla f\!\bigl(w_{i k}\bigr)\bigr\|^{2}
\;\le\;
\frac{2\bigl(f(w_{0})-f^{\star}\bigr)}{\eta\,M}
\;}
\tag{5}
\]
Consequently,
\[
\min_{1\le i\le M}\bigl\|\nabla f\!\bigl(w_{i k}\bigr)\bigr\|
\;\le\;
\sqrt{\frac{2\bigl(f(w_{0})-f^{\star}\bigr)}{\eta\,M}}.
\tag{6}
\]
Since $M\approx T/k$, the bound is $O(1/\sqrt{T})$ for the minimal
gradient norm among snapshots, and the averaged model $\bar w_{T}$ inherits
the same order of stationarity.
\end{theorem}

%%%%%%%%%%%%%%%%%%%%%%%%%%%%%%%%%%%%%%%%%%%%%%%%%%%%%%%%%%%%
\section*{Theoretical Properties}
\begin{enumerate}[label=\arabic*.]
\item \textbf{Stability.} The stepsize condition $\eta\le 1/L$ guarantees that each
gradient step is a non‑expansive map in the sense of the smoothness
inequality, preventing divergence even on highly non‑convex landscapes.
\item \textbf{Hyper‑parameter sensitivity.}
\begin{itemize}
\item \emph{Snapshot period $k$.} Larger $k$ yields fewer snapshots, thus a
coarser averaging (potentially higher variance) but lower memory and
computational overhead.
\item \emph{Stepsize $\eta$.} The bound \eqref{5} is inversely proportional to $\eta$;
a larger admissible $\eta$ (up to $1/L$) gives a tighter guarantee.
\end{itemize}
\item \textbf{Comparison to baselines.}
\begin{itemize}
\item \emph{Plain SGD.} For SGD with the same constant stepsize, the classical
bound on the average gradient norm is $O(1/\sqrt{T})$ but holds for \emph{all}
iterates.  SWA‑I achieves the same order while delivering a single
averaged model that typically generalises better (empirically observed in
SWA literature).
\item \emph{Full weight averaging (SWA).} SWA‑I differs by using linear
interpolation during training; the convergence rate remains unchanged,
yet the interpolation can improve the effective learning rate between
snapshots.
\end{itemize}
\end{enumerate}

%%%%%%%%%%%%%%%%%%%%%%%%%%%%%%%%%%%%%%%%%%%%%%%%%%%%%%%%%%%%
\section*{Preliminaries (Definitions and Lemmas)}
\begin{definition}[Gradient descent step]
For any $w\in\mathbb{R}^{d}$, the GD update with stepsize $\eta$ is
\[
\mathcal{G}_{\eta}(w):=w-\eta\nabla f(w).
\tag{7}
\]
\end{definition}

\begin{lemma}[Descent Lemma]\label{lem:descent}
If $f$ is $L$‑smooth and $\eta\le 1/L$, then for any $w$
\[
f\bigl(\mathcal{G}_{\eta}(w)\bigr)\le
f(w)-\frac{\eta}{2}\,\|\nabla f(w)\|^{2}.
\tag{8}
\]
\end{lemma}
\begin{proof}
From smoothness \eqref{A1'} with $x=w$ and $y=\mathcal{G}_{\eta}(w)=w-\eta\nabla f(w)$,
\[
\begin{aligned}
f\bigl(w-\eta\nabla f(w)\bigr)
&\le f(w)-\eta\|\nabla f(w)\|^{2}
+\frac{L}{2}\eta^{2}\|\nabla f(w)\|^{2}\\[2mm]
&=f(w)-\underbrace{\Bigl(\eta-\frac{L\eta^{2}}{2}\Bigr)}_{\ge\eta/2}
\|\nabla f(w)\|^{2} \qquad\text{[since }\eta\le 1/L\text{]}\\[2mm]
&\le f(w)-\frac{\eta}{2}\|\nabla f(w)\|^{2}.
\end{aligned}
\]
\end{proof}

%%%%%%%%%%%%%%%%%%%%%%%%%%%%%%%%%%%%%%%%%%%%%%%%%%%%%%%%%%%%
\section*{Main Proof (Step‑by‑Step Derivation)}
We prove Theorem \ref{thm:main}.  All expectations are omitted because we
consider full‑batch gradients; the same steps hold for unbiased stochastic
gradients after taking expectations.

\noindent\textbf{Notation.}
Let $t_{i}=i k$ be the $i$‑th snapshot index.
Define $D_{i}:=f(w_{t_{i-1}})-f(w_{t_{i}})$ for $i\ge1$.

\begin{enumerate}
\item[Step 1.] Apply Lemma \ref{lem:descent} to every GD step
between two successive snapshots.  For $t\in\{t_{i-1},\dots,t_{i}-1\}$,
\[
f(w_{t+1})\le f(w_{t})-\frac{\eta}{2}\|\nabla f(w_{t})\|^{2}.
\tag{9}
\]
\item[Step 2.] Sum \eqref{9} over the $k$ steps that constitute one snapshot
interval:
\[
\begin{aligned}
f(w_{t_{i}}) &\le f(w_{t_{i-1}})
-\frac{\eta}{2}\sum_{t=t_{i-1}}^{t_{i}-1}\|\nabla f(w_{t})\|^{2}\\[1mm]
\Longrightarrow\quad
D_{i}&:=f(w_{t_{i-1}})-f(w_{t_{i}})
\ge\frac{\eta}{2}\sum_{t=t_{i-1}}^{t_{i}-1}\|\nabla f(w_{t})\|^{2}.
\tag{10}
\end{aligned}
\]
\item[Step 3.] Because $f$ is bounded below by $f^{\star}$,
\[
\sum_{i=1}^{M} D_{i}=f(w_{0})-f(w_{t_{M}})
\le f(w_{0})-f^{\star}.
\tag{11}
\]
\item[Step 4.] Combine \eqref{10} and \eqref{11}:
\[
\frac{\eta}{2}\sum_{i=1}^{M}\sum_{t=t_{i-1}}^{t_{i}-1}
\|\nabla f(w_{t})\|^{2}
\le f(w_{0})-f^{\star}.
\tag{12}
\]
\item[Step 5.] The left‑hand side of \eqref{12} contains $k$ gradient norms
per snapshot.  Discard all but the last one in each block (the gradient at
the snapshot point $w_{t_{i}}$) to obtain a valid lower bound:
\[
\sum_{i=1}^{M}\|\nabla f(w_{t_{i}})\|^{2}
\;\le\;
\sum_{i=1}^{M}\sum_{t=t_{i-1}}^{t_{i}-1}\|\nabla f(w_{t})\|^{2}.
\tag{13}
\]
\item[Step 6.] Insert \eqref{13} into \eqref{12} and divide by $M$:
\[
\frac{1}{M}\sum_{i=1}^{M}\|\nabla f(w_{t_{i}})\|^{2}
\le
\frac{2\bigl(f(w_{0})-f^{\star}\bigr)}{\eta\,M}.
\tag{14}
\]
Equation \eqref{14} is exactly the bound \eqref{5} claimed in the theorem.
\item[Step 7.] Taking the square root of the right‑hand side of \eqref{14}
and applying the elementary inequality $\min_{i}a_{i}\le\frac{1}{M}\sum_{i}a_{i}$
gives the minimal‑gradient guarantee \eqref{6}.
\end{enumerate}
Thus every step follows directly from the smoothness assumption and
the descent lemma, completing the proof. \hfill $\square$

%%%%%%%%%%%%%%%%%%%%%%%%%%%%%%%%%%%%%%%%%%%%%%%%%%%%%%%%%%%%
\section*{Rate Derivation (With Constants)}
From \eqref{5} and $M\approx T/k$ we have
\[
\frac{1}{M}\sum_{i=1}^{M}\|\nabla f(w_{i k})\|^{2}
\le
\frac{2k\bigl(f(w_{0})-f^{\star}\bigr)}{\eta\,T}
\;=\;
\mathcal{O}\!\left(\frac{1}{T}\right).
\]
Consequently,
\[
\min_{0\le t\le T}\|\nabla f(w_{t})\|
\;\le\;
\mathcal{O}\!\left(\frac{1}{\sqrt{T}}\right),
\]
which matches the optimal rate for first‑order methods on smooth
non‑convex objectives.

%%%%%%%%%%%%%%%%%%%%%%%%%%%%%%%%%%%%%%%%%%%%%%%%%%%%%%%%%%%%
\section*{Special Cases}
\begin{itemize}
\item \textbf{Convex $L$‑smooth $f$.}  Since every snapshot is a feasible point,
the averaging step \eqref{4} yields a point $\bar w_{T}$ that satisfies the
standard $O(1/T)$ function‑value convergence of SGD with constant stepsize.
\item \textbf{Deterministic full‑batch gradients.}  The proof simplifies
because expectations disappear; the bound \eqref{5} holds deterministically.
\item \textbf{Stochastic gradients with bounded variance $\sigma^{2}$.}
  Adding $\mathbb{E}\|g_{t}-\nabla f(w_{t})\|^{2}\le\sigma^{2}$ to the
  analysis introduces an additional term $\frac{\eta\sigma^{2}}{2}$ on the
  right‑hand side of \eqref{12}, leading to the familiar
  $O\bigl(\frac{1}{\sqrt{T}}+\eta\sigma^{2}\bigr)$ rate.
\end{itemize}

%%%%%%%%%%%%%%%%%%%%%%%%%%%%%%%%%%%%%%%%%%%%%%%%%%%%%%%%%%%%
\section*{Discussion}
The convergence theorem shows that, despite the lack of convexity,
SWA‑I guarantees that at least one snapshot attains a gradient norm that
decays as $O(1/\sqrt{T})$.  The final averaged model $\bar w_{T}$ inherits
this stationarity property because the average of vectors with bounded
norm cannot be larger than the maximum norm.  The linear interpolation
\eqref{3} does not affect the asymptotic rate; it merely provides a
smooth trajectory for inference between snapshots and can be interpreted
as a form of implicit learning‑rate scheduling.

\end{document}